This chapter summarizes the progress made so far, identifies the challenges that
remain, and lays out a concrete plan --- with work packages, a timeline, and
expected contributions --- for completing the research.

\section{Summary of Progress}
\label{sec:plan:progress}
\addcontentsline{tocheb}{section}{\protect\numberline{\secnumforhebrewtoc}{סיכום ההתקדמות}}

The preliminary work has established that interpretable, server-side risk scoring
for MCP is feasible. We have (i) formalized the server-as-victim threat model
(\Cref{sec:framework:threat_model}); (ii) designed the factor-based scoring model
and its atomic-operation taxonomy (\Cref{sec:framework:model,sec:framework:atomic});
(iii) implemented the full static and request-time pipeline driven by a single
deterministic local model; and (iv) run an initial multi-oracle evaluation
(\Cref{chapter:experiments}). The results are encouraging where the problem is
well-defined --- $97\%$ agreement with reviewed judgments on tool impact and a
quadratic-weighted $\kappa$ of $+0.66$ against a human panel, second among eight
frameworks --- and, importantly, show that content-aware scoring separates live
attacks where capability-only scores such as CVSS do not. This is sufficient
evidence that the direction is worth completing.

\section{Remaining Challenges}
\label{sec:plan:challenges}
\addcontentsline{tocheb}{section}{\protect\numberline{\secnumforhebrewtoc}{אתגרים שנותרו}}

The preliminary evaluation also exposed the problems the remaining research must
solve. The most significant is the \emph{fragility of the composite band}:
$85\%$ of scored cells change band under a single $\pm1$ perturbation of their
inputs, concentrated at the \textit{medium} threshold, which undermines a gate
that must act on the band. Second, the request-time evaluation rests on
\emph{weak, intent-based labels} (AUC $0.64$), and the static model
systematically under-weights read-only exfiltration. Third, calls are scored
\emph{individually}, so cumulative risk across a sequence of low-risk actions is
not yet modeled. Fourth, reliance on a \emph{single local model} for both scoring
and part of the evaluation introduces a circularity we mitigated but did not
eliminate. These four challenges define the work packages below.

\section{Planned Work}
\label{sec:plan:work}
\addcontentsline{tocheb}{section}{\protect\numberline{\secnumforhebrewtoc}{עבודה מתוכננת}}

\textbf{WP1 --- Stabilize the banding formula.} Replace the hard integer
cut-points with calibrated, monotone soft thresholds, or learn the band
boundaries against the human consensus, and re-run the sensitivity analysis to
target a substantial reduction in the $85\%$ fragility.

\textbf{WP2 --- A contextual request-time scorer.} Build the dynamic scorer that
reads the actual arguments and recent call history to close the
read-only-exfiltration gap, operationalizing the likelihood sub-factors of
\Cref{sec:framework:misuse} from real runtime logs rather than fixed constants.
The external-benchmark result suggests this is where the discriminative power
lies.

\textbf{WP3 --- Multi-step chain modeling.} Extend the scorer to accumulate risk
across a session, so that reconnaissance followed by a single exfiltration is
scored as the campaign it is rather than as isolated low-risk calls.

\textbf{WP4 --- Strengthened, multi-model evaluation.} Replace the weak dynamic
labels with an audited ground truth, broaden the oracle panel, and repeat the
study across multiple model families to remove the single-model circularity and
test generalization.

\textbf{WP5 --- Deployment study.} Integrate the scorer as a gate in front of a
live MCP server and measure end-to-end latency and overhead, addressing the
practical constraints of cost and speed that a real deployment imposes.

\section{Timeline}
\label{sec:plan:timeline}
\addcontentsline{tocheb}{section}{\protect\numberline{\secnumforhebrewtoc}{לוח זמנים}}

We plan the remaining work over four phases, as summarized in
\Cref{tab:timeline}. WP1 and WP2 are prioritized because they address the two
weaknesses that most limit a deployable gate; WP3--WP5 build on their outputs.

\begin{table}[H]
    \centering
    \caption{Planned schedule for the remaining research.}
    \label{tab:timeline}
    \begin{tabular}{l l}
        \hline
        \textbf{Phase} & \textbf{Focus} \\
        \hline
        Months 1--3   & WP1: formula stabilization and re-analysis \\
        Months 3--7   & WP2: contextual request-time scorer \\
        Months 6--9   & WP3: multi-step chain modeling \\
        Months 8--11  & WP4: strengthened, multi-model evaluation \\
        Months 10--12 & WP5: deployment study and thesis writing \\
        \hline
    \end{tabular}
\end{table}

\section{Expected Contributions}
\label{sec:plan:contributions}
\addcontentsline{tocheb}{section}{\protect\numberline{\secnumforhebrewtoc}{תרומות צפויות}}

On completion, the research is expected to contribute: (i) a server-centric
threat model and a factor-based, interpretable risk-scoring model for MCP tool
interactions; (ii) a stabilized banding formula with quantified robustness; (iii)
a contextual request-time scorer that closes the exfiltration gap and models
multi-step campaigns; (iv) a strengthened, multi-model evaluation against human
and framework oracles; and (v) a deployable proof-of-concept gate with measured
overhead. Together these would provide the first graduated, interpretable,
server-side risk-scoring framework for the Model Context Protocol, giving servers
a computable signal to gate, throttle, or deny risky agent requests before they
execute.
